%======================================================================
% Nithin Jambula - Resume: Robotics Software Engineer (2 pages)
% Compile with pdfLaTeX on Overleaf. Single column, ATS friendly.
%======================================================================
\documentclass[a4paper,10pt]{article}

\usepackage[a4paper,left=0.6in,right=0.6in,top=0.5in,bottom=0.5in]{geometry}
\usepackage[T1]{fontenc}
\usepackage{lmodern}
\usepackage{xcolor}
\usepackage{titlesec}
\usepackage{enumitem}
\usepackage[hidelinks]{hyperref}
\input{glyphtounicode}
\pdfgentounicode=1   % lets ATS software read the PDF text correctly

\definecolor{accent}{RGB}{31,78,163}
\hypersetup{colorlinks=true, urlcolor=accent, linkcolor=accent,
  pdftitle={Nithin Jambula - Robotics Software Engineer}, pdfauthor={Nithin Jambula},
  pdfkeywords={Robotics Software Engineer, ROS 2, MuJoCo, Webots, Behavior Trees, Sim-to-Real, Perception, AprilTag, Python, C++}}

\pagestyle{empty}
\setlength{\parindent}{0pt}
\raggedright

\titleformat{\section}{\color{accent}\large\bfseries}{}{0em}{}[{\color{accent}\titlerule[0.6pt]}]
\titlespacing*{\section}{0pt}{9pt}{4pt}
\setlist[itemize]{leftmargin=1.3em, itemsep=1pt, topsep=2pt, parsep=0pt, partopsep=0pt}

% \job{Organisation}{Dates}{Role}{Location}
\newcommand{\job}[4]{\textbf{#1}\hfill #2\\\textit{#3}\hfill\textit{#4}\par\vspace{-2pt}}
% \project{Name}{Dates}{Tools and links}
\newcommand{\project}[3]{\textbf{#1}\hfill #2\\{\small #3}\par\vspace{-2pt}}
\newcommand{\gap}{\vspace{5pt}}

\begin{document}

%----------------------------- HEADER ---------------------------------
\begin{center}
  {\color{accent}\Huge\textbf{Nithin Jambula}}\\[3pt]
  {\large Robotics Software Engineer}\\[4pt]
  \small
  +91 93476 32259 \,$|$\,
  \href{mailto:nithinjambula89@gmail.com}{nithinjambula89@gmail.com} \,$|$\,
  \href{https://linkedin.com/in/nithin-jambula}{linkedin.com/in/nithin-jambula} \,$|$\,
  \href{https://github.com/nithin434}{github.com/nithin434} \,$|$\,
  \href{https://nithin434.github.io}{nithin434.github.io}
\end{center}

%----------------------------- SUMMARY --------------------------------
\section{Summary}
I'm a software engineer who loves writing the software that lets robots see, move and make decisions on their own.
I enjoy working on robot behavior, perception and the testing tools that make robots dependable, and I'm always curious about how a robot's software stays clean and easy to change as it grows.
I like simple, well-structured code that is easy to test, easy to debug and easy to trust.

%---------------------------- EXPERIENCE ------------------------------
\section{Experience}

\job{Ananta Technologies}{Dec 2025 -- Present}{Software Engineer -- Robotics \& Autonomy}{India}
\begin{itemize}
  \item Own the camera-guided robot positioning module (AprilTag pose estimation and motion sequencing) as its main maintainer, keeping it accurate and reliable across multiple autonomous pilot programs.
  \item Designed behavior trees and decision logic that turn high-level mission plans into clear, step-by-step robot actions, with retries and fallbacks so robots recover from errors instead of stopping.
  \item Built and expanded a MuJoCo and Webots simulation test suite to 50+ scenarios covering motion, navigation and failure recovery, cutting validation time by 35\%.
  \item Automated the test runs so new robot behaviors are checked in simulation before they reach hardware, catching problems early and making releases safer.
  \item Took 2 robots from simulation to the production floor, handling on-robot tuning, sim-to-real fixes and post-launch debugging.
  \item Rebuilt the communication layer for older CNC machines and H2D printers that ran closed software, and connected them to the in-house automation pipeline, removing manual operator steps on the print floor.
\end{itemize}
\gap

\job{Vulcan: Self-Driving Campus EV (VIT-AP Initiative)}{Jan 2023 -- Present}{Software Lead -- Perception \& Navigation}{Amaravati, India}
\begin{itemize}
  \item Lead the student team building the software for a self-driving electric campus vehicle: perception, lane following and driving decisions.
  \item Trained a road-scene segmentation model (DeepLabV3+ on Mapillary) to 67\% accuracy and tuned it to run within the limited compute of an NVIDIA Jetson.
  \item Built lane-following and turn-handling control tuned for campus roads, using a lightweight behavior-tree controller.
  \item Trained driving policies in the CARLA simulator with reinforcement learning (SAC) and refined them with imitation learning (DAgger) for smoother lane keeping and obstacle avoidance.
  \item Combined camera, LiDAR and radar data for obstacle awareness, and showcased the working vehicle live at VLaunchpad 2025.
\end{itemize}
\gap

\job{AIR Centre, VIT-AP University}{Jun 2024 -- Nov 2025}{Student Research Engineer -- Applied AI}{Amaravati, India}
\begin{itemize}
  \item Built AI agents that plan, act and learn from feedback to test software systems on their own, combining LLM reasoning with rule-based checks.
  \item Found 34\% more issues than standard tools and cut manual testing effort by 70\%; published the work at IEEE AIRC 2026.
  \item Built and deployed a real-time AI voice calling agent (speech-to-text, LLM replies, text-to-speech) on Asterisk and ViciDial, and demoed it end to end.
\end{itemize}

%----------------------------- PROJECTS -------------------------------
\section{Projects}

\project{AprilTag Docking Rover (ROS 2)}{Dec 2025}{Python, ROS 2, AprilTag, OpenCV, rosbag \,$|$\, \href{https://github.com/nithin434/apriltag_rover}{GitHub}}
\begin{itemize}
  \item Built a ROS 2 rover that finds and docks with its station on its own, using real-time AprilTag detection and pose estimation.
  \item Wrote a multi-stage mission controller (state machine) for approach, alignment and docking, with each step configurable without code changes.
  \item Added simulation worlds, rosbag recording and tools that turn each run into a demo video, making it easy to review and debug missions.
\end{itemize}
\gap

\project{EchoSight: Navigation Glasses for the Visually Impaired (Patent Pending)}{Dec 2024 -- Present}{YOLOv8, Raspberry Pi, OpenCV, Roboflow, GCP, Text-to-Speech \,$|$\, \href{https://drive.google.com/file/d/1r_dGlaLp-_ueKEicPEHzzcnC8IBGrJQT/view?usp=sharing}{Demo}}
\begin{itemize}
  \item Built wearable glasses that detect obstacles and give spoken directions through Bluetooth earbuds.
  \item Fine-tuned YOLOv8 on 50+ custom object classes and optimised it to run at 45+ FPS on a Raspberry Pi.
  \item Added distance and direction estimation from a single camera, with location-aware voice alerts that update as the user moves.
  \item Awardee at the Vikas 2024 Innovation Challenge, showcased at inter-college innovation competitions, and filed for a patent.
\end{itemize}
\gap

\project{16fps: Multi-Agent Video Generation Platform}{Oct 2025 -- Present}{Python, LangChain, Stable Video Diffusion, AnimateDiff, GCP}
\begin{itemize}
  \item Built a team of AI agents that plan, use tools and check their own work to turn a text prompt into a consistent video over a minute long.
  \item Renders multiple videos in parallel on zero-downtime cloud infrastructure; sold commercial rights and access to a media production company.
\end{itemize}
\gap

\project{ScreenAutomate: Gesture and Voice Control}{Oct 2023 -- Jan 2024}{Python, OpenCV, MediaPipe, Speech Recognition \,$|$\, \href{https://github.com/nithin434/Screen-Automate}{GitHub}}
\begin{itemize}
  \item Built a hands-free way to control a computer with hand gestures tracked by camera and spoken commands, designed for accessibility.
  \item Added voice-driven web search, system monitoring and screen automation in real time.
\end{itemize}
\gap

\project{Face Recognition System}{May 2025}{Python, OpenCV DNN, OpenFace, scikit-learn \,$|$\, \href{https://github.com/nithin434/face_recog_AP}{GitHub}}
\begin{itemize}
  \item Built a face recognition pipeline that detects faces with a deep learning detector, turns each face into an embedding, and identifies people with a trained classifier.
  \item Works on both still images and live video, and can learn new people by retraining on a few photos.
\end{itemize}
\gap

\project{DeepShield: Deepfake Detection}{Jul 2023 -- Nov 2023}{Python, PyTorch, CNN, GAN, OpenCV}
\begin{itemize}
  \item Built a CNN video model that detects face-swap deepfakes, improving accuracy by 22\% and cutting false positives by 18\%, with a synthetic video generator that added 40\% more training variety.
  \item Published at IEEE iSAI-NLP 2025.
\end{itemize}

%------------------------------ SKILLS --------------------------------
\section{Skills}
\begin{itemize}
  \item \textbf{Languages:} Python, C++, Bash, SQL
  \item \textbf{Robot Software:} ROS 2, Behavior Trees, State Machines, Mission Planning, Motion Sequencing, Visual Servoing, AprilTag Pose Estimation
  \item \textbf{Simulation \& Testing:} MuJoCo, Webots, CARLA, Simulation Test Automation, Sim-to-Real Transfer, rosbag
  \item \textbf{Perception:} OpenCV, YOLOv8, DeepLabV3+, Semantic Segmentation, Monocular Depth Estimation, Sensor Fusion (Camera, LiDAR, Radar)
  \item \textbf{Planning \& Learning:} Path Planning, Control Systems, Reinforcement Learning (SAC), Imitation Learning (DAgger), PyTorch, TensorFlow, ONNX
  \item \textbf{Hardware \& Tools:} NVIDIA Jetson, Raspberry Pi, CNC and 3D Printer Automation, Linux, Docker, Git, CI/CD
\end{itemize}

%------------------- PUBLICATIONS, PATENT, AWARDS ---------------------
\section{Publications, Patent \& Awards}
\begin{itemize}
  \item \textbf{N. Jambula}, S. C. Sethuraman, et al., ``A Neuro-Symbolic AI Framework for Adaptive and Explainable API Security Testing,'' \textit{7th International Conference on Artificial Intelligence, Robotics, and Control (AIRC)}, IEEE, 2026.
        DOI: \href{https://doi.org/10.1109/AIRC69745.2026.11631384}{10.1109/AIRC69745.2026.11631384}
  \item \textbf{N. Jambula}, N. K. Reddy, K. Pandey, M. Mahesh, ``Securing Reality: AI-Driven Detection of Deepfakes and Counterfeit Currency,'' \textit{20th International Joint Symposium on Artificial Intelligence and Natural Language Processing (iSAI-NLP)}, IEEE, 2025.
        DOI: \href{https://doi.org/10.1109/ISAI-NLP66160.2025.11320726}{10.1109/ISAI-NLP66160.2025.11320726}
  \item \textbf{Patent Pending:} EchoSight, an AI-powered wearable navigation system for the visually impaired.
  \item \textbf{Awardee}, Vikas 2024 Innovation Challenge (EchoSight) \,$|$\, \textbf{Winner}, Engineering Clinics 2024 (real-time threat detection framework)
\end{itemize}

%----------------------- CERTIFICATIONS \& LEADERSHIP ------------------
\section{Certifications \& Leadership}
\begin{itemize}
  \item \href{https://drive.google.com/file/d/1bYBUAIvUOMfzw3lgQpruA0vFzStBH2Kx/view?usp=sharing}{Machine Learning Specialization} -- Stanford \& DeepLearning.AI \,$|$\, \href{https://www.cert.devtown.in/verify/Z1GAbEE}{Object Detection Bootcamp} -- Google
  \item IBM: \href{https://courses.cognitiveclass.ai/certificates/b24fd91a07c14d8ba2b6dbf0d6a6c6c5}{Deep Learning Fundamentals}, \href{https://www.credly.com/badges/91a4bb99-9a17-4c7c-b000-3c7e54c148eb/print}{Machine Learning with Python}, \href{https://www.credly.com/badges/1d70a2bb-7afa-4ef8-9e2f-a945321c0e51/print}{Machine Learning Specialist -- Advanced}
  \item \textbf{President}, \href{https://acsvitap.in/}{ACS International Student Chapter}, VIT-AP \hfill Aug 2023 -- Present
  \item \textbf{Marketing Member}, \href{https://sedsvitap.in/index.html}{SEDS VIT-AP} \hfill Jan 2023 -- Present
\end{itemize}

%---------------------------- EDUCATION -------------------------------
\section{Education}
\textbf{VIT-AP University}, Amaravati \hfill 2023 -- 2027\\
B.Tech in Computer Science and Engineering \hfill CGPA: 8.0 / 10

\end{document}
